Your final course grade is determined by the following components:

\begin{center}
\begin{tabular}{ll}
\toprule
Component & Weight \\
\midrule
Midterm Exam (May 22)            & 20\% \\
Final Exam (June 5)              & 25\% \\
Data Projects (2)                & 20\% (10\% each) \\
Homework (lowest score dropped)  & 25\% \\
Participation (highest 11 of 14) & 10\% \\
\bottomrule
\end{tabular}
\end{center}

Letter grades will be assigned according to the following minimum-grade thresholds:

\begin{center}
\begin{tabular}{*{10}{c}}
\toprule
A & A$-$ & B$+$ & B & B$-$ & C$+$ & C & C$-$ & D & F \\
\midrule
94 & 90 & 87 & 84 & 80 & 77 & 74 & 70 & 60 & 0 \\
\bottomrule
\end{tabular}
\end{center}

Final grades \textbf{will not be rounded}. Thus, a final grade of 93.99 will earn a letter grade of A$-$.

\subsection*{Exams}
Two 90-minute, paper-based exams will be administered to assess your mastery of the application of economic theory and the empirical facts presented in this course. The midterm exam will be assessed during the normal class time on \textbf{Wednesday, May 22, 9:30 a.m.--11:00 a.m. in Ivester E101}. The final exam will also be during the normal class time on \textbf{Wednesday, June 5, 9:30 a.m.--11:00 a.m. in Ivester E101}.

If you have a known scheduling conflict with either exam date, please notify the instructor as soon as possible.

Make-up exams can be scheduled for those who face a sudden illness or emergency on the day of the exam or participate in a pre-scheduled, university-sanctioned event (e.g., varsity athletics). Proper documentation of your illness, emergency, or scheduling conflict must be presented to the instructor before the make-up exam can be taken. If there are compelling privacy concerns for you to not disclose the cause of your absence to the instructor, you may coordinate with the office of Student Care and Outreach to provide such documentation.

Make-up exams will not be offered to students who miss an exam for reasons other than those stated above, and they will earn a zero on that exam.

Students with disabilities or who otherwise need accommodation to complete the exams must register with the Disability Resource Center (DRC) \textbf{prior to the exam}---accommodations cannot be applied retroactively, and I will not make accommodations for students who are not registered with the DRC. Accommodations may take several days to process, so I recommend contacting the DRC as soon as possible. See Accommodation for Disabilities in this syllabus for more information.

\subsection*{Data Projects}
Economics is a social science, meaning that economic theory is tested scientifically using empirical observation and data. Further, an understanding of the macroeconomy requires learning key empirical facts that motivate economists, policymakers, and business leaders. Students will complete 2 data projects during the semester. Completing these projects requires only the use of Microsoft Word and Excel, or Google Docs and Sheets, all of which can be accessed for free via the University.

Tentatively, the due dates for each project are (1) Friday, May 24 and (2) Tuesday, June 4. Each data project will be posted to eLC one week prior to its due date. \textbf{Late submissions will not be accepted.} Detailed instructions related to these projects will be provided during the course.

\subsection*{Homework}
Homework in this course consists of 7 online assignments on the Achieve platform. Homework assignments will assess your understanding of the material covered in lecture and offer immediate feedback on your preparedness for the midterm and final exams.

Homework assignments generally will be due at 11:59 p.m. on the class day \textbf{after} the in-class material is finished. For example, the in-class material for Gains from Trade is scheduled to be completed on Wednesday, May 15, so the corresponding homework will be due at 11:59 p.m. on Thursday, May 16. \textbf{Late submissions will not be accepted.}

You are welcome to work with your classmates on the homework assignments, and you may refer to the textbook or other course materials while you complete these assignments.

I will drop the lowest score of the 7 homework assignments for each student. Scores on the top 6 graded homework assignments will be equally weighted when calculating your final course grade.

\subsection*{Participation}
A participation grade will be recorded each class day (other than the dates of the midterm and final exam) on the following scale:

\begin{ul}
    \item \textbf{3 pts.} Satisfactorily completed the participation assignment.
    \item \textbf{2 pts.} Attempted the participation assignment.
    \item \textbf{0 pts.} Did not attempt the participation assignment.
\end{ul}

Participation assignments will be announced each day and generally consist of an in-class activity that involves practicing the concepts discussed in that day's lecture. Thus, whether you earn 3 points or 0 points for a given day's participation assignment will largely be determined by whether you attended that day's class.

The aggregate participation grade will be calculated by dividing the sum of your top 11 participation scores by 33. There are 14 class meetings at which participation will be recorded, meaning the lowest 3 will be dropped when calculating your participation grade.
